% !TeX program = lualatex
\documentclass[11pt,a4paper]{article}

% ========= حزمة =========
\usepackage[english, bidi=default, layout=counters]{babel}
\usepackage{amsmath,amssymb,amsfonts,amsthm,mathtools}
\usepackage{geometry}
\geometry{left=1.25cm,right=1.25cm,top=1.25cm,bottom=3.1cm}
\usepackage{booktabs}
\usepackage{array}
\usepackage{hyperref}
\usepackage{url}
\usepackage{graphicx}
\usepackage{caption}
\usepackage{subcaption}
\usepackage{float}
\usepackage{siunitx}
\usepackage{bm}
\usepackage{pgfplots}
\pgfplotsset{compat=1.18}
\usepackage{xcolor}
\usepackage{titlesec}
\usepackage{fancyhdr}
\usepackage{microtype}
\usepackage{fontspec}
\usepackage{parskip}

% ========= اللغات =========
\babelprovide[import, main, maparabic, alph=alph]{arabic}
\babelprovide[import, onchar=ids fonts]{english}

% ========= الخطوط =========
\defaultfontfeatures{Ligatures=TeX}
\setmainfont{Amiri}[
Script=Arabic,
Mapping=arabicdigits,
Ligatures=TeX,
BoldFont=Amiri-Bold,
ItalicFont=Amiri-Italic,
BoldItalicFont=Amiri-BoldItalic
]
\setsansfont{Amiri}[
Script=Arabic,
Mapping=arabicdigits
]
\newfontfamily{\arabicfonttt}{Amiri}[
Script=Arabic,
Mapping=arabicdigits
]

% خطوط للنص اللاتيني (الإنجليزية)
\babelfont[english]{rm}{Times New Roman}
\babelfont[english]{sf}{Arial}
\babelfont[english]{tt}{Courier New}

% ========= ألوان =========
\definecolor{skyblue}{RGB}{0,102,204}
\definecolor{darkblue}{RGB}{0,0,139}
\definecolor{gold}{RGB}{255,215,0}

% ========= أوامر YUCT =========
\newcommand{\Keff}{K_{\mathrm{eff}}}
\newcommand{\kappac}{\kappa_c}
\newcommand{\eps}{\varepsilon}
\newcommand{\dint}{D+I\!\cdot\!R}
\newcommand{\YPSDC}{YPSDC}
\newcommand{\YUCT}{YUCT}

% ========= نظريات =========
\newtheorem{theorem}{نظرية}[section]
\newtheorem{definition}[theorem]{تعريف}
\newtheorem{proposition}[theorem]{اقتراح}
\newtheorem{prediction}[theorem]{تنبؤ أعمى}

% ========= تنسيق العناوين =========
\titleformat{\section}{\LARGE\bfseries\color{darkblue}}{\thesection}{1em}{}
\titleformat{\subsection}{\normalsize\bfseries\color{darkblue}}{\thesubsection}{1em}{}
\titleformat{\subsubsection}{\normalsize\itshape}{\thesubsubsection}{1em}{}

% ========= رؤوس وتذييلات =========
\fancypagestyle{firstpage}{%
	\fancyhf{}%
	\renewcommand{\headrulewidth}{-5pt}%
	\renewcommand{\footrulewidth}{-5pt}%
	\fancyfoot[C]{%
		\begin{minipage}[t]{\textwidth}
			\centering
			\textcolor{gold}{\rule{\textwidth}{1.6pt}}\\[5pt]
			{\small \textsuperscript{1}أبحاث ياكوشيف، \url{https://yuct.org/}} \hfill
			{\small بروتوكول YPSDC: \url{https://ypsdc.com/}}\\[-2pt]
			\textcolor{gold}{\rule{\textwidth}{1.6pt}}\\[5pt]
			\textbf{نظرية ياكوشيف الموحدة للتنسيق 2026} \hfill \thepage\\[10pt]
		\end{minipage}%
	}%
}

\fancypagestyle{plain}{%
	\fancyhf{}%
	\renewcommand{\headrulewidth}{0pt}%
	\renewcommand{\footrulewidth}{0pt}%
	\fancyfoot[C]{%
		\begin{minipage}[t]{\textwidth}
			\centering
			\textcolor{gold}{\rule{\textwidth}{1.6pt}}\\[4pt]
			\textbf{نظرية ياكوشيف الموحدة للتنسيق 2026} \hfill \thepage\\[4pt]
		\end{minipage}%
	}%
}

\pagestyle{plain}
\setlength{\parindent}{0pt}
\setlength{\parskip}{1ex}
\raggedbottom

\hypersetup{
	colorlinks=true,
	linkcolor=blue,
	citecolor=blue,
	urlcolor=blue,
}

% ========= رأس المقال (YUCT logo) =========
\newcommand{\makeheader}{%
	\noindent
	\begin{minipage}{0.17\textwidth}
		\raggedright
		{\sffamily\bfseries\color{skyblue}\fontsize{46}{46}\selectfont YUCT}%
	\end{minipage}%
	\hspace{30pt}
	\begin{minipage}{0.32\textwidth}
		\raggedright
		{\sffamily\fontsize{11}{13}\selectfont نظرية ياكوشيف\\ الموحدة\\ للتنسيق}%
	\end{minipage}%
	\hfill
	\begin{minipage}{0.38\textwidth}
		\raggedleft
		{\sffamily\bfseries ملحق U}\\
		{\sffamily مارس 2026}\\
		{\sffamily\footnotesize DOI: \url{https://doi.org/10.5281/zenodo.18444598}}%
	\end{minipage}
	\par\vspace{5pt}
	\noindent\textcolor{gold}{\rule{\textwidth}{1.6pt}}
	\par\vspace{0pt}
}

\begin{document}
	\thispagestyle{firstpage}
	\makeheader
	
	\vspace{0cm}
	{\LARGE\bfseries الاتصالات الثقالية القائمة على كفاءة التنسيق: \\
		تجاوز الحدود الكلاسيكية دون انتهاك السببية \par}
	\vspace{0.2cm}
	
	{\large أليكسي ف. ياكوشيف\footnote{أبحاث ياكوشيف، \url{https://yuct.org/}}} \\
	\vspace{0.0cm}
	
	{\color{darkblue}\textbf{\LARGE المستخلص}} \par
	{\large
		\noindent
		نقترح نهجاً جديداً للاتصالات بعيدة المدى قائماً على نظرية ياكوشيف الموحدة للتنسيق (YUCT). من خلال معاملة حقل الجاذبية كقاموس شامل طبيعي، نُظهر أن أدلة قصيرة (اضطرابات ثقالية) يمكنها تنسيق الحالات بين مراقبين بعيدين بكفاءة تنسيق \(\Keff\) تتجاوز رسمياً حد سرعة الضوء، دون انتهاك السببية. يكمن المفتاح في الفصل بين توزيع القاموس دون اتصال (حقل الجاذبية الموجود مسبقاً) وإرسال الدليل عبر الاتصال (موجة ثقالية بسرعة الضوء). تخترق الموجات الثقالية المادة دون توهين، مما يوفر ميزة كبيرة على الإشارات الكهرومغناطيسية. نشتق المعادلات الأساسية، ونقدر الكفاءات الممكنة، ونقترح اختباراً تجريبياً باستخدام مقاييس جاذبية عالية الدقة. يفتح هذا العمل طريقاً نحو فئة جديدة من أنظمة الاتصالات حيث يتحدد زمن الانتقال بالمعالجة المحلية بدلاً من المسافة.
	}
	
	\vspace{0.1cm}
	\noindent
	{\color{darkblue}\textbf{الكلمات المفتاحية:}} YUCT، اتصالات ثقالية، كفاءة التنسيق، YPSDC، موجات ثقالية، سببية، قاموس، دليل.
	
	\vspace{0.5cm}
	\noindent
	{\small نُسخة مختصرة من هذا العمل نُشرت سابقاً ومتاحة على \\ \url{https://doi.org/10.5281/zenodo.18362308}}
	\vspace{0.5cm}
	
	\noindent
	\textcopyright\ 2026 أبحاث ياكوشيف. جميع الحقوق محفوظة.
	
	\tableofcontents
	\newpage
	
	\section{مقدمة}
	\label{sec:intro}
	
	تعتمد أنظمة الاتصالات الكلاسيكية على الموجات الكهرومغناطيسية (الراديو، الليزر) أو إشارات أخرى تنتشر بسرعة الضوء أو أقل. بالنسبة لمهمات الفضاء العميق، يفرض هذا الحد الأساسي زمن انتقال لا مفر منه: تحتاج الإشارة إلى المريخ ما بين 4 و24 دقيقة، مما يعيق التحكم الفوري ونقل البيانات. حتى مفاهيم مثل الانتقال الآني الكمومي لا تتجاوز هذا القيد، لأنها تتطلب قناة كلاسيكية محدودة بـ \(c\).
	
	تقدم نظرية ياكوشيف الموحدة للتنسيق (YUCT) \cite{Yakushev2026} نقلة نوعية: فهي تميز بين \emph{نقل البيانات} (النقل الفيزيائي للبتات) و\emph{تنسيق الحالات} (محاذاة المعرفة الموجودة مسبقاً). الفكرة الأساسية هي \textbf{بروتوكول \(\YPSDC\) ذو المرحلتين}: مرحلة غير متصلة توزع قاموساً مشتركاً \(D\)؛ ومرحلة متصلة تنقل فقط دليلاً قصيراً \(\kappa\) يُنشط المدخلة المقابلة في \(D\). كفاءة التنسيق
	\begin{equation}
		\Keff = \frac{H(\text{الفعل المُنشط})}{H(\text{الدليل المرسل})}
	\end{equation}
	يمكن أن تكون كبيرة بشكل اعتباطي لأن القاموس يحمل معظم المعلومات. المهم أن الدليل نفسه ينتشر بسرعة لا تتجاوز \(c\)؛ التنسيق الفوري الظاهري ينشأ من القاموس الموجود مسبقاً.
	
	في هذا المقال نستكشف إمكانية استخدام \emph{حقل الجاذبية} كقاموس طبيعي شامل. كل جسم ذي كتلة يشارك في التفاعل الثقالي الكوني، والذي يمكن رؤيته كقاموس موزع باستمرار. بترميز المعلومات في اضطرابات صغيرة لحقل الجاذبية (مثلاً بتعديل كتلة)، يمكننا إرسال أدلة تُفك ترميزها عند المستقبِل باستخدام نفس القاموس الثقالي. نظراً لأن الموجات الثقالية تتفاعل بشكل ضعيف للغاية مع المادة، فهي تخترق العوائق دون توهين — ميزة حاسمة على الإشارات الكهرومغناطيسية. علاوة على ذلك، يربط حقل الجاذبية جميع الأجسام قبلياً، لذا يمكن لكفاءة التنسيق، من حيث المبدأ، أن تتجاوز حاجز سرعة الضوء دون انتهاك السببية.
	
	\section{أسس YUCT}
	\label{sec:foundations}
	
	\subsection{بروتوكول \(\YPSDC\)}
	\label{subsec:ypsdc}
	
	ليشترك مراقبان \(A\) و\(B\) في قاموس \(\mathcal{D} = \{ (\kappa_i, A_i) \}\) يربط الأدلة \(\kappa_i\) بالأفعال \(A_i\). يُوزع القاموس دون اتصال (ربما بسرعة دون ضوئية). أثناء الاتصال المتصل، يرسل \(A\) الدليل \(\kappa_i\) عبر قناة فيزيائية؛ ويستقبله \(B\) ويسترجع فوراً \(A_i\) من \(\mathcal{D}\). الزمن الكلي للتنسيق هو زمن إرسال \(\kappa_i\) مضافاً إليه زمن البحث، والذي يمكن جعله مهملاً.
	
	\subsection{كفاءة التنسيق}
	\label{subsec:keff}
	
	إذا كان كل فعل \(A_i\) يحتوي على \(n\) بت من المعلومات والدليل \(\kappa_i\) بطول \(\ell\) بت، فإن \(\Keff = n/\ell\). بشكل عام،
	\begin{equation}
		\Keff = \frac{H(\mathcal{A})}{H(\mathcal{K})},
	\end{equation}
	حيث \(H\) هي إنتروبيا شانون. لقاموس مضغوط بأمثلية، يمكن أن تصل \(\Keff\) إلى قيم \(10^6\) أو أكثر (مثلاً في بروتوكولات الاتصالات الحديثة).
	
	\subsection{السببية و \(K_{\mathrm{eff}} > 1\)}
	\label{subsec:causality}
	
	قد يوحي تفسير ساذج لـ \(\Keff > 1\) بنقل معلومات أسرع من الضوء. لكن \(\Keff\) يقيس نسبة المعلومات المُنشطة إلى البتات المرسلة، وليس سرعة البتات نفسها. ينتقل الدليل \(\kappa\) بسرعة \(v \le c\)؛ والمعلومات الإضافية تأتي من القاموس الموزع مسبقاً. تُحفظ السببية لأن القاموس أُنشئ سابقاً بوسائل دون ضوئية. وهكذا يمكن أن تكون \(\Keff\) كبيرة اعتباطياً دون انتهاك النسبية الخاصة.
	
	\section{الجاذبية كقاموس شامل}
	\label{sec:gravity}
	
	\subsection{لماذا الجاذبية؟}
	
	لحقل الجاذبية خواص فريدة تجعله مرشحاً مثالياً لقاموس شامل:
	\begin{itemize}
		\item \textbf{الشمولية}: كل جسم ذي كتلة-طاقة يشارك في الجاذبية. الحقل موجود في كل مكان.
		\item \textbf{الاختراق}: تتفاعل الموجات الثقالية بشكل ضعيف جداً مع المادة بحيث يمكنها عبور الكواكب والنجوم والمجرات دون توهين أو تشتت يُذكر.
		\item \textbf{البنية الموجودة مسبقاً}: يُشفر حقل الجاذبية أصلاً توزيع الكتل؛ ويمكن اعتباره قاموساً تشترك فيه جميع الأجسام منذ لحظة نشوئها.
	\end{itemize}
	
	\subsection{حقل الجاذبية كقاموس}
	
	في YUCT، القاموس هو مجموعة من الحالات الممكنة. للاتصالات الثقالية، نُفسر تنسور المترية \(g_{\mu\nu}\) (أو اضطراباته) كالقاموس. يعرف المستقبِل حقل الجاذبية المحيط (مثلاً حقل الأرض، حقل الشمس، إلخ) ويمكنه كشف الانحرافات عنه. يُعدل المرسِل حقل الجاذبية المحلي — مثلاً بتحريك كتلة، تغيير شكلها، أو إثارة رنين — مولِّداً بذلك اضطراباً صغيراً \(\delta g_{\mu\nu}\). يعمل هذا الاضطراب كدليل \(\kappa\). المستقبِل، المزود بمقياس جاذبية حساس أو مقياس تداخل، يقيس \(\delta g_{\mu\nu}\) ويستخدم القاموس المشترك (حقل الخلفية المعروف) ليفسره كفعل محدد.
	
	\subsection{تشبيه \(\YPSDC\) في الجاذبية}
	
	\begin{itemize}
		\item \textbf{المرحلة غير المتصلة}: يتأسس حقل الجاذبية الخلفي بتوزيع جميع الكتل في الكون. هذا الحقل "موزع" تلقائياً منذ بداية الزمن.
		\item \textbf{المرحلة المتصلة}: يُنشئ المرسِل اضطراباً محلياً (دليلاً) ينتشر نحو الخارج بسرعة الضوء (كموجة ثقالية). المستقبِل، عند كشف الاضطراب، يقارنه بالخلفية المعروفة ويستخلص المعلومات المُشفرة.
	\end{itemize}
	
	نظراً لأن الاضطراب ينتقل بسرعة \(c\)، لا يوجد انتشار إشارة فائق للضوء. لكن \emph{التنسيق} بين المرسِل والمستقبِل — أي حقيقة أن اضطراباً صغيراً ينقل كمية كبيرة محتملة من المعلومات — يمكن أن يكون فعالاً للغاية، أي \(\Keff \gg 1\).
	
	\section{نموذج رياضي للاتصالات الثقالية}
	\label{sec:model}
	
	\subsection{ترميز الإشارة}
	
	لنفترض أن المرسِل يُعدل كتلة محلية \(M(t)\) وفق شيفرة متفق عليها مسبقاً. للتبسيط، باعتبار ترميز ثنائي: تغير \(\Delta M\) لمدة \(\tau\) يُمثل "1"، بينما عدم التغير يُمثل "0". سعة الموجة الثقالية الناتجة على مسافة \(r\) تُعطى تقريباً بـ
	\begin{equation}
		h(t) \sim \frac{G}{c^4 r} \frac{d^2}{dt^2} \left[ M(t) \right],
	\end{equation}
	حيث \(G\) ثابت نيوتن. تدفق الطاقة ضئيل جداً، لكن الكواشف الحديثة (LIGO، مراصد فضائية مستقبلية) يمكنها قياس انفعالات حتى \(10^{-22}\).
	
	\subsection{سعة القناة وكفاءة التنسيق}
	
	سعة قناة الموجات الثقالية محدودة بأرضية الضجيج للكاشف وعرض النطاق المتاح. لكن \emph{سعة التنسيق} أعلى لأن كل شكل موجة مُستقبل يمكن مطابقته مع مكتبة من القوالب الممكنة. ليحتو القاموس على \(N\) رسالة ممكنة، كل منها موصوفة بقالب شكل موجة فريد. يقوم المستقبِل بترابط متبادل للإشارة الواردة مع جميع القوالب؛ الدليل هو أساساً معرف القالب. إذا صُممت القوالب لتكون متعامدة، فإن عدد البتات المنقولة لكل حدث مُستقبل هو \(\log_2 N\). الطاقة أو الزمن اللازم لإرسال الموجة الفيزيائية مستقل عن \(N\). وبالتالي،
	\begin{equation}
		\Keff = \frac{\log_2 N}{\text{(طول البت الفعال للإشارة الفيزيائية)}}.
	\end{equation}
	مثلاً، إذا كان لدينا \(10^6\) قالب وكل إشارة فيزيائية تشغل ثانية واحدة من البيانات بمعدل عينات 1 kHz، فإن حجم البيانات الخام هو \(1000\) بت، لكن المعلومات المنقولة هي \(\log_2 10^6 \approx 20\) بت. وهكذا \(\Keff \approx 20/1000 = 0.02\)، وهو أقل من 1. لتحقيق \(\Keff > 1\)، نحتاج إلى قوالب قصيرة جداً مقارنة بما تحمله من معلومات. من حيث المبدأ، يمكن استخدام هوائيات ثقالية رنانة مضبوطة على ترددات محددة؛ عندها تُنشط نبضة قصيرة عند تردد معين مدخلة قاموسية محددة، ويكون الدليل أساساً هو التردد نفسه. نظراً لأن التردد يمكن قياسه بدقة عالية، يمكن أن يكون عدد الترددات المميزة هائلاً، مما يعطي \(\Keff\) كبيراً.
	
	\subsection{الضجيج وقياس الخطأ}
	
	وفقاً لقانون الخطأ الشامل في YUCT \cite{AppendixL}،
	\begin{equation}
		\eps = \kappac \alpha \Keff^{-\beta}, \quad \beta = 0.67,
		\label{eq:universal-scaling}
	\end{equation}
	حيث \(\eps\) هو احتمال سوء التعرف. يضع هذا حداً أساسياً على مقدار \(\Keff\) الذي يمكن بلوغه قبل أن تصبح الأخطاء غير مقبولة. بالنسبة للاتصالات الثقالية، سيسيطر على معدل الخطأ ضجيج الكاشف وعدم تطابق القوالب، لكن القانون الشامل يوحي بوجود \(\Keff\) أمثل.
	
	\section{مقارنة مع الاتصالات الكهرومغناطيسية}
	\label{sec:comparison}
	
	\subsection{التوهين والاختراق}
	
	تُحجب الموجات الكهرومغناطيسية أو تُشتت بالمادة: يمكن أن تُمتص الموجات الراديوية بالأيونوسفيرات أو الماء أو الأجزاء الداخلية للكواكب؛ وتحتاج الإشارات الضوئية إلى خط رؤية. تعبر الموجات الثقالية أي شيء دون تأثر يُذكر. هذا يعني أن وصلة ثقالية يمكن أن تعمل حتى عندما يكون المرسِل والمستقبِل على جانبين متقابلين من كوكب أو نجم أو عائق آخر.
	
	\subsection{متطلبات الطاقة}
	
	يتطلب توليد موجات ثقالية قابلة للكشف كتلاً وتسارعات هائلة. لا يمكن للتكنولوجيا الحالية إنتاج موجات ثقالية اصطناعية قوية بما يكفي للاتصالات عبر المسافات الفلكية. لكن بالنسبة للمسافات بين الكواكب داخل المجموعة الشمسية، قد تكون الانفعالات المطلوبة في متناول أنظمة ميكانيكية عالية الطاقة مستقبلية (مثلاً كتل دوارة كبيرة أو مذبذبات مدفوعة نووياً). علاوة على ذلك، إذا استخدمنا مصادر طبيعية موجودة (مثل النجوم النابضة) كمنارات، فهي تُوفر أصلاً نوعاً من "الحامل" الثقالي الذي يمكن تعديله.
	
	\subsection{زمن الانتقال}
	
	تنتشر كل من الموجات الكهرومغناطيسية والثقالية بسرعة \(c\)، لذا فإن زمن الانتقال أحادي الاتجاه متطابق. ميزة القناة الثقالية ليست في السرعة بل في \textbf{الكفاءة}: لأن القاموس موجود مسبقاً، يمكن لدليل قصير أن ينقل رسالة معقدة، ضاغطاً المعلومات فعلياً. هذا لا يقلل زمن الانتقال، لكنه يقلل عرض النطاق والطاقة اللازمين لكل بت.
	
	\subsection{إمكانات كفاءة التنسيق}
	
	بالنسبة للاتصالات الراديوية، أقصى \(\Keff\) محدود بسعة القناة وتعقيد التعديل. مع الموجات الثقالية، يمكن أن تكون \(\Keff\) المحتملة أعلى بكثير لأن القاموس (حقل الجاذبية الخلفي) غني ومستقر للغاية. مثلاً، يمكن أن يعمل الكمون الثقالي لنظام الأرض-قمر كقاموس بملايين الحالات المميزة (تشكيلات مدارية مختلفة). تعديل مدار قمر صناعي صغير يمكن أن يُشفر معلومات يستطيع مراقب بعيد، يعرف التقويم الفلكي، فك ترميزها بكفاءة عالية.
	
	\subsection{لا تماثلية القناة الثقالية}
	\label{subsec:asymmetry}
	
	يكمن فارق جوهري بين الاتصالات الكهرومغناطيسية والثقالية في لا تماثلية متطلبات المرسِل والمستقبِل. بينما تتطلب وصلة الليزر توجيهاً دقيقاً وكاشفاً ضوئياً محاذياً للحزمة الواردة، يعمل كاشف الموجات الثقالية — مثل مقياس التداخل الليزري — بشكل شامل الاتجاهات. إنه يراقب انفعال الزمكان باستمرار وينتظر إشارة مميزة ("الشيفرة") دون الحاجة لمعرفة اتجاه المصدر \cite{LIGO}. هذا يجعل الاستقبال الثقالي أبسط جوهرياً من حيث التوجيه والتتبع.
	
	لكن توليد موجة ثقالية قابلة للكشف يتطلب طاقة هائلة وإزاحات كتلية ضخمة. يمكن للمحطات الأرضية، من حيث المبدأ، استضافة كتل رنانة كبيرة أو معدّلات مدارية (مثل مقاييس تداخل بمقياس كيلومتر مثل LIGO، أو مراصد فضائية مستقبلية). أما مركبة فضائية صغيرة، فتفتقر إلى الكتلة والقدرة اللازمتين لإنتاج انفعال يُكشف على مسافات بين كوكبية. وبالتالي فإن الوصلة غير متماثلة بطبيعتها: من مرسِل أرضي قوي (أو مداري كبير) إلى مستقبِل سلبي يصغي فقط. يتطلب الاتصال ثنائي الاتجاه إما مرسِلاً ضخماً مماثلاً على المركبة الفضائية — غير ممكن حالياً — أو نوعاً جديداً من بواعث الموجات الثقالية، مثل تجويف رنان عالي التردد بطاقة نووية، وهو ما يبقى تخمينياً.
	
	لا يقلل هذا التماثل من قيمة القناة في تطبيقات عديدة: يمكن لمسابير الفضاء العميق تلقي أوامر معقدة (أدلة) دون الحاجة إلى مرسِل قوي على متنها، بينما يمكن إرسال القياس عن بعد بالراديو التقليدي. وهكذا تعمل الوصلة الثقالية الهابطة كقناة بث عالية الكفاءة، مكملة الوصلات الصاعدة الموجودة.
	
	قد تمكن التطورات المستقبلية في الأنظمة الميكانيكية عالية الطاقة، أو استخدام مصادر موجات ثقالية طبيعية (مثل النجوم النابضة المعدّلة) كمنارات، من إنشاء وصلات متماثلة في النهاية. في الوقت الحالي، يجب أن ينصب التركيز على إثبات جدوى الاتصال الثقالي أحادي الاتجاه وقياس كفاءة تنسيقه \(K_{\mathrm{eff}}\) في تجربة مضبوطة.
	
	\section{دور كثافة حقل التنسيق في انتشار الإشارة والديناميك الثقالي}
	\label{sec:coord-density}
	
	\subsection{كثافة حقل التنسيق كمُعدِّل لهندسة الزمكان}
	
	في إطار YUCT، الجاذبية ليست مجرد خاصية هندسية لزمكان فارغ، بل تجلٍ لحقل التنسيق الأساسي \(\Psi_{MN}\). تلعب كثافة هذا الحقل، المُشار إليها بـ \(\rho_K = \langle \Psi_{MN}\Psi^{MN} \rangle\)، دوراً مماثلاً لقرينة انكسار لتدفق المعلومات. في مناطق كثافة التنسيق العالية — كتلك القريبة من أجسام ضخمة، أو في أنظمة كمومية مترابطة، أو أثناء التحولات الطورية في الكون المبكر — يمكن أن تتعدل المترية الفعالة التي تحكم انتشار الأدلة.
	
	انطلاقاً من لاغرانجيان كامل ذي 19 بعداً (الملحق A) وحدود ترابطه (مثلاً \(\mathcal{L}_{329}\) الذي يصف الترابط الرنيني بين القطاعات)، يمكننا اشتقاق فعل فعال رباعي الأبعاد حيث يُساهم حقل التنسيق في مترية منبثقة:
	\begin{equation}
		g_{\mu\nu}^{\mathrm{eff}} = g_{\mu\nu} + \alpha \, \rho_K \, T_{\mu\nu}^{\mathrm{coord}},
	\end{equation}
	حيث \(\alpha\) ثابت لا بعدي و\(T_{\mu\nu}^{\mathrm{coord}}\) تنسور الطاقة-الاندفاع لحقل التنسيق. ونتيجة لذلك، فإن سرعة الدليل (اضطراب صغير في الحقل) لم تعد تماماً \(c\) بل تصبح
	\begin{equation}
		v_{\mathrm{index}} = c \left(1 + \beta \, \rho_K + \cdots \right),
	\end{equation}
	حيث \(\beta\) ثابت آخر يرتبط بقوة الترابط \(\kappa_{0,103}\) بين قطاعي الجاذبية والمعلومات. هذا التأثير صغير للغاية في الظروف العادية لكنه قد يصبح قابلاً للقياس في السياقات الفيزيائية الفلكية أو الكونية القصوى.
	
	\subsection{التعزيز الرنيني ووجود قنوات مفضلة}
	
	وجود حدود لاخطية في اللاغرانجيان، مثل حد \(-\epsilon \Phi^3\) في \(\mathcal{L}_{329}\)، يسمح بتضخيم وسيطي للإشارات عند ترددات محددة \(\omega_{mn}\). في مناطق حيث \(\rho_K\) كبيرة، يمكن أن تُنشئ هذه الرنينات "أنهاراً" من التنسيق المعزز — اتجاهات أو ترددات مفضلة ينتشر عبرها الدليل بأقل فقدان وحتى بتضخيم، مستمداً الطاقة من حقل التنسيق نفسه. هذا مماثل لوسط ربح ليزري، لكن التضخيم هنا يحدث في النسيج الهندسي للزمكان.
	
	بالنسبة لوصلة اتصال ثقالي، يعني هذا أن الكفاءة \(\Keff\) ليست محددة فقط بحجم القاموس، بل أيضاً بكثافة التنسيق المحلية. إذا استطاع مرسِل تعديل كتلة عند تردد رنين \(\omega_{\mathrm{res}}\) يطابق تردداً ذاتياً لحقل التنسيق على طول المسار إلى المستقبِل، فسيتم توجيه الإشارة تفضيلياً، مما يقلل بشكل كبير من الاستطاعة المطلوبة عند المصدر. المستقبِل، المضبوط على نفس التردد، يعمل ككاشف رنيني، مستخلصاً الدليل من ضجيج الخلفية.
	
	\subsection{تنبؤات لانتشار الإشارة}
	
	من الاعتبارات أعلاه، نحصل على ثلاث تنبؤات قابلة للاختبار تُكمل ما ورد في القسم \ref{sec:comparison}:
	
	\begin{prediction}[اعتماد اتجاهي لسرعة الموجات الثقالية]
		إذا كانت كثافة حقل التنسيق \(\rho_K\) غير موحدة الخواص (مثلاً بسبب بنية واسعة النطاق أو إطار مفضل)، فإن سرعة الموجات الثقالية (وبالتالي أي دليل ثقالي) يجب أن تُظهر اعتماداً اتجاهياً صغيراً. ربط أزمنة وصول إشارات الموجات الثقالية من اتجاهات مختلفة في السماء يمكن أن يكشف عن تباين من مرتبة \(\Delta v/c \sim 10^{-18}\)–\(10^{-20}\)، في متناول شبكات مستقبلية من كواشف الجيل الثالث (مرصد أينشتاين، Cosmic Explorer).
	\end{prediction}
	
	\begin{prediction}[ترددات رنينية في نطاق الميلي هرتز]
		يتنبأ الحد اللاخطي \(\epsilon \Phi^3\) في \(\mathcal{L}_{329}\) بوجود ترددات رنينية ضيقة في نطاق الميلي هرتز (حيث قد يكون لحقل التنسيق للبنى المجرية أنماط ذاتية). باستخدام تجاويف بصرية فائقة الاستقرار كما اقترح في \cite{Barontini2025}، يجب أن يكشف البحث عن مثل هذه الرنينات في الترابط المتبادل لكاشفين بعيدين عن قمم عند ترددات \(\omega_{mn}\) لا تُعزى لأي مصدر فلكي معروف.
	\end{prediction}
	
	\begin{prediction}[تضخيم الإشارات الاصطناعية بحقل التنسيق]
		إذا تم ضبط مصدر موجات ثقالية اصطناعي (مثلاً كتلة دوارة كبيرة) على تردد رنين لحقل التنسيق المحلي (مثلاً نظام الأرض-قمر)، فيجب أن تتجاوز شدة الإشارة المستقبلة عند كاشف بعيد تنبؤ صيغة الرباعي القياسية بمعامل \(\sim Q\)، حيث \(Q\) عامل جودة الرنين (يمكن أن يصل إلى \(10^3\)–\(10^6\)). سيشكل هذا برهاناً مباشراً على تضخيم حقل التنسيق.
	\end{prediction}
	
	\subsection{الارتباط بقانون القياس الشامل}
	
	ينطبق قانون قياس الخطأ الشامل \(\eps = \kappac\alpha\Keff^{-\beta}\) (معادلة \ref{eq:universal-scaling}) أيضاً على انتشار الأدلة. في منطقة ذات \(\rho_K\) معززة، تزداد \(\Keff\) الفعالة لطول دليل معطى، مما يقلل احتمال الخطأ. يُوفر هذا آلية طبيعية لتحقيق كفاءات تنسيق عالية للغاية دون انتهاك السببية: القاموس موجود مسبقاً، وينتقل الدليل عبر وسط يدعم توصيله بنشاط.
	
	\section{اختبار تجريبي مقترح}
	\label{sec:experiment}
	
	\subsection{الإعداد}
	
	وضع مقياسي جاذبية عاليي الحساسية (مثلاً مقاييس جاذبية فائقة التوصيل أو مقاييس تداخل ذرية) في موقعين بعيدين على الأرض، لنقل على بعد 1000 كم. في أحد الموقعين، تعديل كتلة كبيرة (مثلاً رقاص بعدة أطنان) بنمط محدد مسبقاً. تنتشر الإشارة الثقالية بسرعة \(c\) وتُكشف في الموقع الآخر. بمقارنة الإشارة المستقبلة مع النمط المعروف، يمكننا قياس كفاءة التنسيق الفعالة ومقارنتها بالتنبؤ النظري.
	
	\subsection{نتائج متوقعة}
	
	نتوقع أنه حتى بقاموس بسيط نسبياً (مثلاً مجموعة رموز سعة)، ستتجاوز \(\Keff\) الحد الكلاسيكي لنفس الإنفاق الطاقي. ستختبر التجربة أيضاً قانون الخطأ الشامل: كلما زدنا عدد الرموز (زادت \(\Keff\))، يجب أن يتبع معدل الخطأ \(\eps \propto \Keff^{-0.67}\).
	
	\subsection{تحديات}
	
	التحدي الرئيسي هو السعة الضئيلة للموجات الثقالية الاصطناعية. على مسافة 1000 كم، تنتج كتلة 1000 طن تتذبذب بسعة 1 م عند 1 هرتز انفعالاً من مرتبة \(10^{-23}\)، عند حد القابلية الحالية للكشف. قد تحقق كواشف الموجات الثقالية المستقبلية (مثلاً مرصد أينشتاين) الحساسية المطلوبة.
	
	\section{مناقشة}
	\label{sec:discussion}
	
	يقدم مفهوم الاتصالات الثقالية القائم على YUCT خروجاً جذرياً عن التفكير التقليدي. باستغلال حقل الجاذبية الموجود مسبقاً كقاموس، يمكننا تحقيق كفاءات تنسيق تبدو متجاوزة لسرعة الضوء، مع احترام صارم للسببية. يواجه التطبيق العملي عقبات تقنية هائلة، لكن الإطار النظري سليم وقابل للاختبار.
	
	إذا نجح، يمكن لمثل هذا النظام أن يُمكن اتصالات فورية (من منظور المستخدم) عبر المجموعة الشمسية، لأن زمن الانتقال سيهيمن عليه زمن توليد وكشف الدليل، وليس المسافة. علاوة على ذلك، ولأن الموجات الثقالية تخترق كل مادة، سيكون الاتصال ممكناً من أي مكان، بما في ذلك الأعماق تحت الأرض أو تحت المحيطات.
	
	\section{خلاصة}
	\label{sec:conclusion}
	
	قدمنا إطاراً نظرياً للاتصالات الثقالية قائماً على نظرية ياكوشيف الموحدة للتنسيق. بمعاملة حقل الجاذبية كقاموس شامل طبيعي، يمكن لأدلة قصيرة (اضطرابات ثقالية) تنسيق أفعال بين مراقبين بعيدين بكفاءة \(\Keff\) غير محدودة بسرعة الضوء. الدليل نفسه لا يزال ينتشر بسرعة \(c\)، حافظاً على السببية. تُقدم الموجات الثقالية مزايا فريدة من حيث الاختراق والشمولية، مما يجعلها جذابة لأنظمة اتصالات الفضاء العميق المستقبلية. أوجزنا اختباراً تجريبياً محتملاً باستخدام التكنولوجيا الحالية. يفتح هذا العمل اتجاهاً جديداً في كل من فيزياء الجاذبية ونظرية المعلومات، بآثار عميقة على مستقبل الاتصالات بين النجوم.
	
	\section*{شكر وتقدير}
	
	يشكر المؤلف المشاركين في ندوة YUCT على المناقشات الملهمة.
	
	\begin{thebibliography}{99}
		
		\bibitem{Yakushev2026}
		A. V. Yakushev, \emph{Yakushev Unified Coordination Theory (YUCT)}, Zenodo, \url{https://doi.org/10.5281/zenodo.18444598} (2026).
		
		\bibitem{AppendixL}
		A. V. Yakushev, \emph{Appendix L: Fractal Coordination Error Scaling — A Universal Law from DNA to Cosmology}, Zenodo (2026).
		
		\bibitem{AppendixA}
		A. V. Yakushev, \emph{Appendix A: Practical Guide to Using YUCT V35.0 for Interdisciplinary Modeling and Predictions}, Zenodo (2026).
		
		\bibitem{LLR-IfE}
		J. M. et al., "Lunar Laser Ranging: A Continuing Legacy of the Apollo Program", \emph{Science} \textbf{265}, 482–490 (1994).
		
		\bibitem{LIGO}
		B. P. Abbott et al. (LIGO Scientific Collaboration), "Observation of Gravitational Waves from a Binary Black Hole Merger", \emph{Phys. Rev. Lett.} \textbf{116}, 061102 (2016).
		
		\bibitem{Barontini2025}
		F. Barontini et al., "Detecting millihertz gravitational waves with optical cavities", \emph{Phys. Rev. Lett.} (to be published) (2025).
		
	\end{thebibliography}
	
\end{document}